%\frontmatter
\thispagestyle{empty}
%\chapter*{Agradecimientos}
{\LARGE \textcolor{myblue}{Acerca de la portada}}\\

\begingroup
\setlength{\parindent}{0pt} % Disable indentation within this group

La portada frontal intenta recrear el estado actual del conocimiento sobre las integrinas, combinándolo con un hobby personal: los cubos de Rubik. Las integrinas descubrieron relativamente hace poco tiempo y la información acerca de los mecanismos moleculares y termodinámicos involucrados en la actividad biológica de cada integrina es aún un área en estudio. Incluso, la estructura secundaria y terciaria de ciertos miembros de la gran familia de integrinas no se conoce con exactitud.

La cara derecha del cubo muestra el modelo de la estructura secundaria de las integrinas junto con ciertos ligandos de la matriz extracelular. En la cara inferior se representan algunos ligandos que interactúan con los segmentos citoplasmáticos de las integrinas, además de dos switches, uno de encendido y otro de apagado, que corresponderían al estado activo e inactivo de algunas integrinas. En la cara izquierda, se observa a varios científicos y científicas investigando sobre el tema, para luego plasmar esa información en las aristas, esquinas y centros del cubo. \\\vspace{10cm}


 \tiny {La ilustración tiene un carácter meramente pedagógico y simplificado. No pretende representar con exactitud la realidad de los procesos biológicos de las integrinas.}


\endgroup